% !TeX root = ../complete.tex
\section{از توان لحظه‌ای تا انرژی یک گذار}
\label{sec:power-energy}

\subsection{توان لحظه‌ای، انرژی و توان متوسط}
\label{sec:power-average}

انرژی با ژول اندازه‌گیری می‌شود. توان، نرخ مبادله انرژی برحسب وات است.
برای نمونه، یک وات یعنی مبادله یک ژول در یک ثانیه. اگر جریان
\(i(t)\) از پایانه مثبت یک عنصر وارد شود و اختلاف پتانسیل دو پایانه
\(v(t)\) باشد، توان جذب‌شده عنصر در زمان \(t\) از حاصل‌ضرب این دو
به دست می‌آید \cite[بخش \secnum{۵.۱.۱}، صفحه ۱۸۲]{weste}.
\begin{equation}
  p(t)=v(t)i(t).
  \label{eq:power-instantaneous}
\end{equation}
در این قرارداد، \(p(t)>0\) یعنی عنصر انرژی می‌گیرد و \(p(t)<0\)
یعنی به مدار انرژی می‌دهد. در تحلیل تغذیه، \(I_{DD}(t)\) را در جهت
خروج از پایانه مثبت منبع تعریف می‌کنیم. بنابراین \(V_{DD}I_{DD}(t)\)
توان تحویلی منبع است. واحد حاصل‌ضرب ولت و آمپر همان وات است.

برای محاسبه انرژی، بازه بین \(t_0\) و \(t_1\) را به تکه‌های زمانی
کوتاه تقسیم کنید. در هر تکه با طول \(\Delta t\)، اگر توان تقریباً
ثابت باشد، انرژی آن تکه تقریباً \(p(t)\Delta t\) است. جمع این
انرژی‌ها در حد تکه‌های بسیار کوتاه همان انتگرال توان است. توان متوسط
توان ثابتی است که در همان مدت، همان انرژی را مبادله می‌کند.
با \(T=t_1-t_0>0\)، از رابطه‌های انتگرال و میانگین استفاده می‌کنیم
\cite[صفحه ۱۸۲]{weste}.
\begin{equation}
  E=\int_{t_0}^{t_1}p(t)\,\mathrm{d}t,
  \qquad
  P_{\mathrm{avg}}=\frac{E}{T}.
  \label{eq:power-interval}
\end{equation}
روی نمودار توان برحسب زمان، انتگرال برابر مساحت علامت‌دار بین منحنی
و محور زمان است. اگر توان در نیمه نخست یک بازه‌ی
\qty{4}{\nano\second} برابر \qty{3}{\milli\watt} و در نیمه دوم
صفر باشد، تنها نیمه نخست در انرژی سهم دارد. تبدیل پیشوندها در حاصل‌ضرب
انرژی به‌صورت زیر دیده می‌شود.
\[
  E=(\qty{3e-3}{\watt})(\qty{2e-9}{\second})
    =\qty{6e-12}{\joule}
    =\qty{6}{\pico\joule}.
\]
این انرژی بر کل زمان \qty{4}{\nano\second} تقسیم می‌شود و توان
متوسط \qty{1.5}{\milli\watt} به دست می‌آید. اگر زمان انتظار بدون
مصرف طولانی‌تر شود، انرژی همان رویداد ثابت می‌ماند و توان متوسط در
بازه‌ی طولانی‌تر کم‌تر می‌شود.

\subsection{انرژی خازن و انرژی تحویلی منبع}
\label{sec:power-capacitor-energy}

یک \lr{Inverter} خروجی ورودی خود را معکوس می‌کند. وقتی ورودی از
یک به صفر می‌رود، \lr{pMOS} خروجی را به تغذیه وصل می‌کند. وقتی
ورودی از صفر به یک می‌رود، \lr{nMOS} خروجی را به زمین وصل می‌کند.
در مدل شکل \ref{fig:power-charge-paths}، بار خروجی خازن ثابت \(C_L\)
است. فرض می‌کنیم در هر گذار فقط مسیر نشان‌داده‌شده هدایت می‌کند،
منبع در \(V_{DD}\) ثابت می‌ماند و خروجی پیش از گذار بعدی به سطح
نهایی خود می‌رسد. این مدل با تقریب بارِ
\NotesTopic{sec:delay-inverter-transient}{پاسخ گذرای \lr{Inverter}}
سازگار است \cite[شکل \secnum{۵.۴}، صفحه ۱۸۳]{weste}.

\begin{figure}[htbp]
  \centering
  \begin{latin}
  \begin{tikzpicture}[x=1cm,y=1cm]
    \foreach \xx/\stage in {0/\(p\mathrm{MOS}\),6/\(n\mathrm{MOS}\)} {
      \begin{scope}[xshift=\xx cm]
        \draw[thick] (0,1.1) -- (0,0) -- (3,0) -- (3,0.46);
        \draw[thick] (2.73,0.46) -- (3.27,0.46);
        \draw[thick] (2.73,0.64) -- (3.27,0.64);
        \draw[thick] (3,0.64) -- (3,1.65) -- (1.55,1.65);
        \draw[thick] (1.55,1.4) rectangle (0.45,1.9);
        \node at (1,1.65) {\stage};
        \draw[thick] (0.45,1.65) -- (0,1.65) -- (0,1.1);
        \draw[thick] (2.72,0) -- (3.28,0);
        \draw[thick] (2.81,-0.12) -- (3.19,-0.12);
        \draw[thick] (2.91,-0.24) -- (3.09,-0.24);
        \fill (3,1.65) circle (1.5pt);
        \node[above] at (3,1.65) {\(v_o\)};
        \node[right] at (3.32,0.55) {\(C_L\)};
      \end{scope}
    }
    \draw[thick,fill=white] (0,0.75) circle (0.35);
    \node at (0,0.87) {\scriptsize \(+\)};
    \node at (0,0.63) {\scriptsize \(-\)};
    \node[left] at (-0.4,0.75) {\(V_{DD}\)};
    \draw[->,blue!70!black,thick] (1.65,1.39) -- (2.65,1.39)
      node[midway,below] {\(i_C\)};
    \draw[->,blue!70!black,thick] (8.65,1.39) -- (7.65,1.39)
      node[midway,below] {\(-i_C\)};
    \node[above] at (1.5,2.05) {\rl{شارژ}};
    \node[above] at (7.5,2.05) {\rl{تخلیه}};
  \end{tikzpicture}
  \end{latin}
  \caption{مسیر معادل شارژ و تخلیه بار خروجی یک \lr{Inverter}.
    \(i_C\) در جهت ورود به صفحه بالایی \(C_L\) تعریف شده است.
    در تخلیه، جهت جریان معکوس و بزرگی آن برابر \(-i_C\) است.
    مدار از تحلیل شکل \secnum{۵.۴} کتاب اقتباس شده است \cite{weste}.}
  \label{fig:power-charge-paths}
\end{figure}

طبق \NotesTopic{sec:foundation-capacitor}{مدل خازن ثابت}، بار صفحه
بالایی خازن در ولتاژ \(v_o\) برابر \(q=C_Lv_o\) است و جریان
\(i_C=\mathrm{d}q/\mathrm{d}t\) نرخ تغییر این بار را نشان می‌دهد.
وقتی بسته‌ی کوچک بار \(\mathrm{d}q\) وارد خازن می‌شود، از رابطه
\eqref{eq:power-instantaneous} انرژی جذب‌شده برابر
\(v_o\,\mathrm{d}q\) است. ولتاژ در طول شارژ تغییر می‌کند. بنابراین
برای محاسبه انرژی باید سهم هر بسته بار را با ولتاژ همان لحظه جمع کنیم.
با جای‌گذاری \(\mathrm{d}q=C_L\,\mathrm{d}v_o\)، انرژی ذخیره‌شده
پس از شارژ از صفر تا \(V_{DD}\) چنین است
\cite[رابطه‌های \secnum{۵.۶} و \secnum{۵.۷}، صفحات ۱۸۲ و ۱۸۳]{weste}.
\begin{equation}
  E_C=\int_{0}^{V_{DD}} C_Lv_o\,\mathrm{d}v_o
     =C_L\left[\frac{v_o^2}{2}\right]_{0}^{V_{DD}}
     =\frac{1}{2}C_LV_{DD}^2.
  \label{eq:power-capacitor-energy}
\end{equation}
مشتق \(v_o^2/2\) نسبت به \(v_o\) برابر \(v_o\) است. به همین دلیل
انتگرال با تفریق مقدار \(v_o^2/2\) در دو کران به دست می‌آید. روی نمودار
\(v_o\) برحسب \(q\) نیز همین انرژی مساحت مثلثی به قاعده
\(C_LV_{DD}\) و ارتفاع \(V_{DD}\) است.

در مسیر شارژ، تمام بارِ واردشده به خازن از پایانه مثبت منبع عبور
می‌کند. منبع برای هر واحد بار، انرژی متناظر با ولتاژ ثابت \(V_{DD}\)
را می‌دهد. پس برای یافتن انرژی تحویلی منبع \(E_{DD}\)، بار کل
\(Q=C_LV_{DD}\) را در \(V_{DD}\) ضرب می‌کنیم
\cite[رابطه \secnum{۵.۸}، صفحه ۱۸۳]{weste}
\cite[اسلایدهای ۳ تا ۵]{powerslides}.
\begin{equation}
  E_{DD}=\int_0^Q V_{DD}\,\mathrm{d}q
        =V_{DD}Q=C_LV_{DD}^2.
  \label{eq:power-source-energy}
\end{equation}
واحد \(C_LV_{DD}^2\) برابر فاراد ضرب‌در مربع ولت است. از آنجا که
فاراد برابر کولن بر ولت و ژول برابر ولت ضرب‌در کولن است، حاصل ژول
می‌شود. اختلاف \(E_{DD}-E_C\) در \lr{pMOS} به گرما تبدیل می‌شود.
در تخلیه، تغذیه از مسیر بار جداست و انرژی ذخیره‌شده در \lr{nMOS}
تلف می‌شود. در پایان چرخه، انرژی خازن به مقدار اولیه بازمی‌گردد و
مجموع انرژی تلف‌شده در دو مسیر برابر \(C_LV_{DD}^2\) است.

\subsection{مثال و کنترل موازنه انرژی}
\label{sec:power-cycle-example}

بار \qty{150}{\femto\farad} را با تغذیه \qty{1.0}{\volt} در
نظر بگیرید. این داده‌ها از مثال شکل موج اسلاید ۶ و شکل
\secnum{۵.۵} کتاب گرفته شده‌اند. محاسبه را با فرض‌های مدل شکل
\ref{fig:power-charge-paths} انجام می‌دهیم
\cite[اسلاید ۶]{powerslides}\cite[صفحات ۱۸۳ و ۱۸۴]{weste}.
\[
  E_{DD}=(\qty{150e-15}{\farad})(\qty{1.0}{\volt})^2
        =\qty{150}{\femto\joule},
  \qquad
  E_C=\frac{E_{DD}}{2}=\qty{75}{\femto\joule}.
\]
در شارژ، \qty{75}{\femto\joule} در \lr{pMOS} تلف و
\qty{75}{\femto\joule} در خازن ذخیره می‌شود. در تخلیه، همین انرژی
ذخیره‌شده در \lr{nMOS} به گرما تبدیل می‌شود. با ثابت بودن ظرفیت و
ولتاژ و کامل بودن شارژ و تخلیه، کندتر شدن گذار مقدار بار جابه‌جاشده و
انرژی این مدل را تغییر نمی‌دهد. شکل لحظه‌ای جریان و مدت گذار می‌توانند
تغییر کنند. بخش \ref{sec:power-short-circuit} سهم جریان مستقیم از
تغذیه به زمین را بررسی می‌کند.
